\documentclass[10pt,a4paper]{amsart}
\usepackage{amsmath,amssymb,amsthm,graphicx}
% Retained preamble for the Stacks Project chapter snapshot in this directory.
% The chapter's own preamble.tex is a copy of the shared Stacks Project preamble; the block
% below is sliced MECHANICALLY, line for line, from it. Nothing was added, removed or
% redefined; the wrapper class, the xr-hyper/hyperref wiring and the externaldocument list
% are omitted because the assembler supplies the document class.
%
% lines taken from preamble.tex: 3, 5, 9, 11, 12, 156, 157, 158, ... (51 lines total)
%
% package coverage restored from the chapter's own preamble (2026-10-13): xy, xy-2cell

\usepackage[all]{xy}
\xyoption{2cell}
\UseAllTwocells

\newenvironment{reference}{\comment}{\endcomment}
\newenvironment{slogan}{\comment}{\endcomment}
\newenvironment{history}{\comment}{\endcomment}
\theoremstyle{plain}
\newtheorem{theorem}[subsection]{Theorem}
\newtheorem{proposition}[subsection]{Proposition}
\newtheorem{lemma}[subsection]{Lemma}
\theoremstyle{definition}
\newtheorem{definition}[subsection]{Definition}
\newtheorem{example}[subsection]{Example}
\newtheorem{exercise}[subsection]{Exercise}
\newtheorem{situation}[subsection]{Situation}
\theoremstyle{remark}
\newtheorem{remark}[subsection]{Remark}
\newtheorem{remarks}[subsection]{Remarks}
\numberwithin{equation}{subsection}
\def\lim{\mathop{\mathrm{lim}}\nolimits}
\def\colim{\mathop{\mathrm{colim}}\nolimits}
\def\Spec{\mathop{\mathrm{Spec}}}
\def\Hom{\mathop{\mathrm{Hom}}\nolimits}
\def\Ext{\mathop{\mathrm{Ext}}\nolimits}
\def\SheafHom{\mathop{\mathcal{H}\!\mathit{om}}\nolimits}
\def\SheafExt{\mathop{\mathcal{E}\!\mathit{xt}}\nolimits}
\def\Sch{\mathit{Sch}}
\def\Mor{\mathop{\mathrm{Mor}}\nolimits}
\def\Ob{\mathop{\mathrm{Ob}}\nolimits}
\def\Sh{\mathop{\mathit{Sh}}\nolimits}
\def\NL{\mathop{N\!L}\nolimits}
\def\CH{\mathop{\mathrm{CH}}\nolimits}
\def\proetale{{pro\text{-}\acute{e}tale}}
\def\etale{{\acute{e}tale}}
\def\QCoh{\mathit{QCoh}}
\def\Ker{\mathop{\mathrm{Ker}}}
\def\Im{\mathop{\mathrm{Im}}}
\def\Coker{\mathop{\mathrm{Coker}}}
\def\Coim{\mathop{\mathrm{Coim}}}
\DeclareMathSymbol{\boxtimes}{\mathbin}{AMSa}{"02}
\def\QCohstack{\mathcal{QC}\!\mathit{oh}}
\def\Cohstack{\mathcal{C}\!\mathit{oh}}
\def\Spacesstack{\mathcal{S}\!\mathit{paces}}
\def\Quotfunctor{\mathrm{Quot}}
\def\Hilbfunctor{\mathrm{Hilb}}
\def\Curvesstack{\mathcal{C}\!\mathit{urves}}
\def\Polarizedstack{\mathcal{P}\!\mathit{olarized}}
\def\Complexesstack{\mathcal{C}\!\mathit{omplexes}}
\def\Pic{\mathop{\mathrm{Pic}}\nolimits}
\def\Picardstack{\mathcal{P}\!\mathit{ic}}
\def\Picardfunctor{\mathrm{Pic}}
\def\Deformationcategory{\mathcal{D}\!\mathit{ef}}
\usepackage{cleveref}
\crefname{originallemma}{original lemma}{original lemmas}
\Crefname{originallemma}{Original lemma}{Original lemmas}
\makeatletter
\newlabel{algebra-lemma-CM-over-regular-flat}{{00R4}{source}{}{sourcepointer-1}{}}
\newlabel{algebra-lemma-CM-over-regular-flat@cref}{{[originallemma][0][]00R4}{[1][0][]source}{}{}{}}
\makeatother
\crefname{originallemma}{original lemma}{original lemmas}
\Crefname{originallemma}{Original lemma}{Original lemmas}
\makeatletter
\newlabel{categories-lemma-yoneda}{{001P}{source}{}{sourcepointer-2}{}}
\newlabel{categories-lemma-yoneda@cref}{{[originallemma][0][]001P}{[1][0][]source}{}{}{}}
\makeatother
\crefname{originallemma}{original lemma}{original lemmas}
\Crefname{originallemma}{Original lemma}{Original lemmas}
\makeatletter
\newlabel{groupoids-lemma-group-scheme-over-field-separated}{{047L}{source}{}{sourcepointer-3}{}}
\newlabel{groupoids-lemma-group-scheme-over-field-separated@cref}{{[originallemma][0][]047L}{[1][0][]source}{}{}{}}
\makeatother
\crefname{originallemma}{original lemma}{original lemmas}
\Crefname{originallemma}{Original lemma}{Original lemmas}
\makeatletter
\newlabel{lemma-open-representable}{{0B9X}{source}{}{sourcepointer-4}{}}
\newlabel{lemma-open-representable@cref}{{[originallemma][0][]0B9X}{[1][0][]source}{}{}{}}
\makeatother
\crefname{originallemma}{original lemma}{original lemmas}
\Crefname{originallemma}{Original lemma}{Original lemmas}
\makeatletter
\newlabel{morphisms-lemma-scheme-theoretic-image-is-proper}{{0AH6}{source}{}{sourcepointer-5}{}}
\newlabel{morphisms-lemma-scheme-theoretic-image-is-proper@cref}{{[originallemma][0][]0AH6}{[1][0][]source}{}{}{}}
\makeatother
\crefname{originallemma}{original lemma}{original lemmas}
\Crefname{originallemma}{Original lemma}{Original lemmas}
\makeatletter
\newlabel{morphisms-lemma-smooth-flat-smooth-fibres}{{01V8}{source}{}{sourcepointer-6}{}}
\newlabel{morphisms-lemma-smooth-flat-smooth-fibres@cref}{{[originallemma][0][]01V8}{[1][0][]source}{}{}{}}
\makeatother
\crefname{originallemma}{original lemma}{original lemmas}
\Crefname{originallemma}{Original lemma}{Original lemmas}
\makeatletter
\newlabel{perfect-lemma-chi-locally-constant-geometric}{{0B9T}{source}{}{sourcepointer-7}{}}
\newlabel{perfect-lemma-chi-locally-constant-geometric@cref}{{[originallemma][0][]0B9T}{[1][0][]source}{}{}{}}
\makeatother
\crefname{originallemma}{original lemma}{original lemmas}
\Crefname{originallemma}{Original lemma}{Original lemmas}
\makeatletter
\newlabel{proposition-hilb-d}{{0B9I}{source}{}{sourcepointer-8}{}}
\newlabel{proposition-hilb-d@cref}{{[originallemma][0][]0B9I}{[1][0][]source}{}{}{}}
\makeatother
\crefname{originallemma}{original lemma}{original lemmas}
\Crefname{originallemma}{Original lemma}{Original lemmas}
\makeatletter
\newlabel{proposition-pic-curve}{{0B9Z}{source}{}{sourcepointer-9}{}}
\newlabel{proposition-pic-curve@cref}{{[originallemma][0][]0B9Z}{[1][0][]source}{}{}{}}
\makeatother
\crefname{originallemma}{original lemma}{original lemmas}
\Crefname{originallemma}{Original lemma}{Original lemmas}
\makeatletter
\newlabel{remark-universal-object-hilb-d}{{0B9J}{source}{}{sourcepointer-10}{}}
\newlabel{remark-universal-object-hilb-d@cref}{{[originallemma][0][]0B9J}{[1][0][]source}{}{}{}}
\makeatother
\crefname{originallemma}{original lemma}{original lemmas}
\Crefname{originallemma}{Original lemma}{Original lemmas}
\makeatletter
\newlabel{varieties-lemma-vanishing-degree-2g-and-1-line-bundle}{{0B90}{source}{}{sourcepointer-11}{}}
\newlabel{varieties-lemma-vanishing-degree-2g-and-1-line-bundle@cref}{{[originallemma][0][]0B90}{[1][0][]source}{}{}{}}
\makeatother
\title{The Picard scheme of a curve decomposes into smooth proper pieces}
\author{The Stacks Project Authors}
\date{stacks-project:pic.tex:local-snapshot-20261012; GNU FDL 1.2 or later}
\begin{document}
\maketitle
\noindent\textit{Editorial note.} Statement span is lines 1412-1432 of sources/stacks-pic/pinned-original-source.tex; proof span is lines 1434-1543, the complete author proof (108 body lines).\par\medskip
\noindent\textit{Editorial note.} Navigation stubs: 11 cited-but-not-included labels are supplied as non-load-bearing pointers with their pinned tags where the snapshot tag index has them.\par\medskip
\noindent\textit{Editorial note.} extra\_packages is empty on purpose: the retained preamble already loads xy with options and a second declaration clashes in the assembler wrapper.\par\medskip
\noindent\textit{Editorial note.} Text is retained verbatim, including the author's own typos and placeholders.\par\medskip
\noindent\textit{Editorial note.} Queue index 0 of queue\_pic.json (pb=4839, stmt 1412-1432, proof 1434-1543).\par\medskip
\noindent\textit{Editorial note.} stacks-pic wave-1 (pi01): fresh adjudication accept (2026-10-14); pool per standing queue hypercovering -> pic (SMC reassigned to Zhen per Chief 08:5x); pic.tex 27 labels/23 with proof judged 21 accept / 2 reject / 4 pb=0 env-dup rows excluded; spans cite-free and reference/slogan/history-free so no patchenv needed. Named omission: Author omits the precise formulation and proof of the isomorphism Z = Proj(Sym(Gamma(X,L)\textasciicircum{}*)) - the inverse construction (section gives effective Cartier divisor) is described in-span; named mechanical check.\par\medskip
\noindent\textit{Allowed standard prerequisites (editorial).}\begin{itemize}
\item Explicit prerequisite boundary. The cited inputs are prerequisites; the argument in the retained span is not replaced by any cited result.
\end{itemize}\medskip
\ifdefined\hypertarget\hypertarget{sourcepointer-1}{}\fi
\noindent\textit{Source reference.} Original lemma 00R4: sources/stacks-pic/pinned-original-source.tex, source lines 1543--1543. Cited by the retained statement or proof but not part of the collected spans; supplied only as a navigation pointer. The argument does not depend on its content.\par\medskip
\ifdefined\hypertarget\hypertarget{sourcepointer-2}{}\fi
\noindent\textit{Source reference.} Original lemma 001P: sources/stacks-pic/pinned-original-source.tex, source lines 1543--1543. Cited by the retained statement or proof but not part of the collected spans; supplied only as a navigation pointer. The argument does not depend on its content.\par\medskip
\ifdefined\hypertarget\hypertarget{sourcepointer-3}{}\fi
\noindent\textit{Source reference.} Original lemma 047L: sources/stacks-pic/pinned-original-source.tex, source lines 1543--1543. Cited by the retained statement or proof but not part of the collected spans; supplied only as a navigation pointer. The argument does not depend on its content.\par\medskip
\ifdefined\hypertarget\hypertarget{sourcepointer-4}{}\fi
\noindent\textit{Source reference.} Original lemma 0B9X: sources/stacks-pic/pinned-original-source.tex, source lines 1214--1219. Cited by the retained statement or proof but not part of the collected spans; supplied only as a navigation pointer. The argument does not depend on its content.\par\medskip
\ifdefined\hypertarget\hypertarget{sourcepointer-5}{}\fi
\noindent\textit{Source reference.} Original lemma 0AH6: sources/stacks-pic/pinned-original-source.tex, source lines 1543--1543. Cited by the retained statement or proof but not part of the collected spans; supplied only as a navigation pointer. The argument does not depend on its content.\par\medskip
\ifdefined\hypertarget\hypertarget{sourcepointer-6}{}\fi
\noindent\textit{Source reference.} Original lemma 01V8: sources/stacks-pic/pinned-original-source.tex, source lines 1543--1543. Cited by the retained statement or proof but not part of the collected spans; supplied only as a navigation pointer. The argument does not depend on its content.\par\medskip
\ifdefined\hypertarget\hypertarget{sourcepointer-7}{}\fi
\noindent\textit{Source reference.} Original lemma 0B9T: sources/stacks-pic/pinned-original-source.tex, source lines 1543--1543. Cited by the retained statement or proof but not part of the collected spans; supplied only as a navigation pointer. The argument does not depend on its content.\par\medskip
\ifdefined\hypertarget\hypertarget{sourcepointer-8}{}\fi
\noindent\textit{Source reference.} Original lemma 0B9I: sources/stacks-pic/pinned-original-source.tex, source lines 849--867. Cited by the retained statement or proof but not part of the collected spans; supplied only as a navigation pointer. The argument does not depend on its content.\par\medskip
\ifdefined\hypertarget\hypertarget{sourcepointer-9}{}\fi
\noindent\textit{Source reference.} Original lemma 0B9Z: sources/stacks-pic/pinned-original-source.tex, source lines 1377--1380. Cited by the retained statement or proof but not part of the collected spans; supplied only as a navigation pointer. The argument does not depend on its content.\par\medskip
\ifdefined\hypertarget\hypertarget{sourcepointer-10}{}\fi
\noindent\textit{Source reference.} Original lemma 0B9J: sources/stacks-pic/pinned-original-source.tex, source lines 912--931. Cited by the retained statement or proof but not part of the collected spans; supplied only as a navigation pointer. The argument does not depend on its content.\par\medskip
\ifdefined\hypertarget\hypertarget{sourcepointer-11}{}\fi
\noindent\textit{Source reference.} Original lemma 0B90: sources/stacks-pic/pinned-original-source.tex, source lines 1543--1543. Cited by the retained statement or proof but not part of the collected spans; supplied only as a navigation pointer. The argument does not depend on its content.\par\medskip
\begin{lemma}
\label{lemma-picard-pieces}
Let $k$ be a separably closed field. Let $X$ be a smooth projective
curve of genus $g$ over $k$.
\begin{enumerate}
\item $\underline{\Picardfunctor}_{X/k}$ is a disjoint union of
$g$-dimensional smooth proper varieties $\underline{\Picardfunctor}^d_{X/k}$,
\item $k$-points of $\underline{\Picardfunctor}^d_{X/k}$
correspond to invertible $\mathcal{O}_X$-modules of degree $d$,
\item $\underline{\Picardfunctor}^0_{X/k}$
is an open and closed subgroup scheme,
\item for $d \geq 0$ there is a canonical morphism
$\gamma_d :
\underline{\Hilbfunctor}^d_{X/k} \to \underline{\Picardfunctor}^d_{X/k}$
\item the morphisms $\gamma_d$
are surjective for $d \geq g$ and smooth for $d \geq 2g - 1$,
\item the morphism
$\underline{\Hilbfunctor}^g_{X/k} \to \underline{\Picardfunctor}^g_{X/k}$
is birational.
\end{enumerate}
\end{lemma}

\begin{proof}
Pick a $k$-rational point $\sigma$ of $X$. Recall that $\Picardfunctor_{X/k}$
is isomorphic to the functor $\Picardfunctor_{X/k, \sigma}$. By
Derived Categories of Schemes, Lemma
\ref{perfect-lemma-chi-locally-constant-geometric}
for every $d \in \mathbf{Z}$ there is an open subfunctor
$$
\Picardfunctor^d_{X/k, \sigma} \subset \Picardfunctor_{X/k, \sigma}
$$
whose value on a scheme $T$ over $k$ consists of those
$\mathcal{L} \in \Picardfunctor_{X/k, \sigma}(T)$ such that
$\chi(X_t, \mathcal{L}_t) = d + 1 - g$ and moreover we have
$$
\Picardfunctor_{X/k, \sigma} =
\coprod\nolimits_{d \in \mathbf{Z}} \Picardfunctor^d_{X/k, \sigma}
$$
as fppf sheaves. It follows that the scheme $\underline{\Picardfunctor}_{X/k}$
(which exists by Proposition \ref{proposition-pic-curve})
has a corresponding decomposition
$$
\underline{\Picardfunctor}_{X/k, \sigma} =
\coprod\nolimits_{d \in \mathbf{Z}} \underline{\Picardfunctor}^d_{X/k, \sigma}
$$
where the points of $\underline{\Picardfunctor}^d_{X/k, \sigma}$ correspond
to isomorphism classes of invertible modules of degree $d$ on $X$.

\medskip\noindent
Fix $d \geq 0$. There is a morphism
$$
\gamma_d :
\underline{\Hilbfunctor}^d_{X/k}
\longrightarrow
\underline{\Picardfunctor}^d_{X/k}
$$
coming from the invertible sheaf $\mathcal{O}(D_{univ})$ on
$\underline{\Hilbfunctor}^d_{X/k} \times_k X$
(Remark \ref{remark-universal-object-hilb-d}) by the Yoneda lemma
(Categories, Lemma \ref{categories-lemma-yoneda}).
Our proof of the representability of the Picard functor of $X/k$
in Proposition \ref{proposition-pic-curve} and
Lemma \ref{lemma-open-representable} shows that $\gamma_g$
induces an open immersion on a nonempty open of
$\underline{\Hilbfunctor}^g_{X/k}$. Moreover, the proof shows
that the translates of this open by $k$-rational points of
the group scheme $\underline{\Picardfunctor}_{X/k}$ define
an open covering. Since
$\underline{\Hilbfunctor}^g_{X/K}$ is smooth of dimension $g$
(Proposition \ref{proposition-hilb-d})
over $k$, we conclude that the
group scheme $\underline{\Picardfunctor}_{X/k}$ is smooth of dimension $g$
over $k$.

\medskip\noindent
By Groupoids, Lemma \ref{groupoids-lemma-group-scheme-over-field-separated}
we see that $\underline{\Picardfunctor}_{X/k}$ is separated.
Hence, for every $d \geq 0$, the image of $\gamma_d$
is a proper variety over $k$
(Morphisms, Lemma \ref{morphisms-lemma-scheme-theoretic-image-is-proper}).

\medskip\noindent
Let $d \geq g$. Then for any field extension $K/k$ and any invertible
$\mathcal{O}_{X_K}$-module $\mathcal{L}$ of degree $d$, we see that
$\chi(X_K, \mathcal{L}) = d + 1 - g > 0$. Hence $\mathcal{L}$ has a
nonzero section and we conclude that $\mathcal{L} = \mathcal{O}_{X_K}(D)$
for some divisor $D \subset X_K$ of degree $d$. It follows that
$\gamma_d$ is surjective.

\medskip\noindent
Combining the facts mentioned above we see that
$\underline{\Picardfunctor}^d_{X/k}$ is proper for $d \geq g$.
This finishes the proof of (2) because now we see that
$\underline{\Picardfunctor}^d_{X/k}$ is proper for $d \geq g$ but
then all $\underline{\Picardfunctor}^d_{X/k}$ are proper by translation.

\medskip\noindent
It remains to prove that $\gamma_d$ is smooth for $d \geq 2g - 1$.
Consider an invertible $\mathcal{O}_X$-module $\mathcal{L}$ of degree
$d$. Then the fibre of the point corresponding to $\mathcal{L}$ is
$$
Z = \{D \subset X \mid \mathcal{O}_X(D) \cong \mathcal{L}\} \subset
\underline{\Hilbfunctor}^d_{X/k}
$$
with its natural scheme structure. Since any isomorphism
$\mathcal{O}_X(D) \to \mathcal{L}$ is well defined up
to multiplying by a nonzero scalar, we see that the canonical
section $1 \in \mathcal{O}_X(D)$ is mapped to a section
$s \in \Gamma(X, \mathcal{L})$ well defined up to multiplication
by a nonzero scalar. In this way we obtain a morphism
$$
Z \longrightarrow
\text{Proj}(\text{Sym}(\Gamma(X, \mathcal{L})^*))
$$
(dual because of our conventions). This morphism is an isomorphism,
because given an section of $\mathcal{L}$ we can take the associated
effective Cartier divisor, in other words we can construct an inverse of
the displayed morphism; we omit the precise formulation and proof.
Since $\dim H^0(X, \mathcal{L}) = d + 1 - g$ for every
$\mathcal{L}$ of degree $d \geq 2g - 1$ by
Varieties, Lemma \ref{varieties-lemma-vanishing-degree-2g-and-1-line-bundle}
we see that $\text{Proj}(\text{Sym}(\Gamma(X, \mathcal{L})^*))
\cong \mathbf{P}^{d - g}_k$.
We conclude that $\dim(Z) = \dim(\mathbf{P}^{d - g}_k) = d - g$.
We conclude that the fibres of the morphism $\gamma_d$ all
have dimension equal to the difference of the dimensions of
$\underline{\Hilbfunctor}^d_{X/k}$ and $\underline{\Picardfunctor}^d_{X/k}$.
It follows that $\gamma_d$ is flat, see
Algebra, Lemma \ref{algebra-lemma-CM-over-regular-flat}.
As moreover the fibres are smooth, we conclude that $\gamma_d$
is smooth by Morphisms, Lemma \ref{morphisms-lemma-smooth-flat-smooth-fibres}.
\end{proof}

\end{document}
